\documentclass[conference]{IEEEtran}

\usepackage{cite}
\usepackage{amsmath,amssymb}
\usepackage{graphicx}
\usepackage{array}
\usepackage{url}

\begin{document}

\title{Process Mining for Quality-Aware Smart-Pallet Logistics}

\author{
\IEEEauthorblockN{Boyan Chakarov}
\IEEEauthorblockA{
University of Twente\\
Enschede, The Netherlands\\
Student number: s2714965 {\hspace{0cm}}
}
\and
\IEEEauthorblockN{Johan Tunc}
\IEEEauthorblockA{
University of Twente\\
Enschede, The Netherlands\\
Student number: s3020010{\hspace{0cm}}
}
\and
\IEEEauthorblockN{Sem de Jong}
\IEEEauthorblockA{
University of Twente\\
Enschede, The Netherlands\\
Student number: s2844427{\hspace{0cm}}
}
}

\maketitle

\begin{abstract}
This report applies process mining to a logistics and manufacturing case study involving smart pallets and quality-depreciating products. The process concerns a dynamic pickup-and-delivery system in which products move from an arrival region to a climate-controlled processing region and finally to climate-controlled trucks. The objective is to identify process variants, bottlenecks, deviations, and transport configurations that influence waiting time, cycle time, quality decay, and disposal. Using a structured process mining methodology inspired by PM2, event logs are prepared from simulation-generated data and analyzed using process discovery, variant analysis, expected-lifecycle conformance analysis, performance analysis, disposal analysis, and vehicle-resource analysis. Three experiment configurations are compared: a random-dispatching baseline, a quality-aware balanced-fleet configuration, and an alternative no-UAV fleet configuration. The results show that Experiment 14 provides the strongest overall configuration: it reduces mean total waiting time from 14.32 to 5.54 time units compared with the baseline, reduces mean cycle time, maintains high lifecycle conformance, and results in no disposed products. Experiment 23 is not recommended because it shows lower conformance, high waiting time, high cycle time, and a 4.95\% disposal rate.
\end{abstract}

\begin{IEEEkeywords}
process mining, logistics, smart pallets, quality decay, event logs, simulation, PM4Py, process improvement
\end{IEEEkeywords}

\section{Introduction and Problem Context}

Quality-sensitive logistics processes require fast and reliable transportation decisions because delays can lead to quality loss or product disposal. This project analyzes a simulated logistics and manufacturing process in which products are transported on smart pallets through three regions. Products arrive in Region 1, are transported to Region 2 for climate-controlled processing, and are then transported through Region 3 to climate-controlled trucks.

The practical problem is that products are subject to quality depreciation over time. Therefore, waiting time, vehicle assignment, travel time, dispatching rules, and fleet composition can directly affect product quality and operational performance. The stakeholder perspective used in this report is that of a logistics or plant manager who wants to reduce quality decay and disposal without making the transport system unnecessarily complex.

The research objective is to use process mining techniques to analyze the smart-pallet logistics process and identify process variants, bottlenecks, deviations, and transport configurations associated with better logistics performance.

The research questions are:

\begin{itemize}
\item RQ1: What is the actual product flow, and what are the dominant process variants?
\item RQ2: To what extent do cases follow the expected product lifecycle?
\item RQ3: How do dispatching rules and fleet configurations influence waiting time, cycle time, quality decay, and disposal?
\item RQ4: Which process improvements can reduce quality-related logistics problems while remaining operationally feasible?
\end{itemize}

\section{Related Work}

Process mining enables the discovery, monitoring, and improvement of processes based on event data \cite{van_der_aalst}. It is useful when the actual process behavior is not fully visible from predefined process descriptions. In logistics and manufacturing, process mining can reveal transport bottlenecks, resource behavior, waiting patterns, process variants, and deviations from expected flows.

This project also relates to simulation-generated event logs. Simulation provides a controlled experimental setting, while process mining helps reveal the actual behavior generated by the simulation. The dataset used in this report was created from an agent-based simulation of a logistics cold-chain case study involving quality-depreciating products and smart pallets \cite{bemthuis2020,dataset}. The event logs contain product identifiers, activities, timestamps, vehicles, and quality decay attributes, making them suitable for process discovery, conformance-style analysis, and performance enhancement.

The project follows a structured process mining approach inspired by PM2 \cite{pm2}. PM4Py is used as the main process mining framework because it supports reproducible process mining analysis in Python \cite{pm4py}. The study is also related to hybrid modeling and simulation, since process mining is applied to data produced by a simulation model \cite{hybrid}.

\section{Methodology}

This project follows a structured process mining workflow inspired by PM2. The workflow consists of project planning, data extraction, preprocessing, process discovery, conformance-style checking, enhancement through performance analysis, and improvement recommendation.

\subsection{Project Planning}

The project is framed from the perspective of a logistics manager. The main goal is to identify which configuration of dispatching logic and fleet composition leads to stable product flow, low waiting time, low cycle time, low disposal, and acceptable quality outcomes.

The selected key performance indicators are total waiting time, waiting time in Region 1, waiting time in Region 3, travel time, total time in system, quality decay, product disposal rate, process variant frequency, and expected lifecycle conformance.

\subsection{Data Extraction}

The dataset contains input files, raw simulation output, warm-up-filtered output, event logs, and summarized experiment results. The main folders used in this project are \texttt{LogFilesProductWarmupFilter}, \texttt{OutputWarmupFilter}, and the experiment settings files. The product-level event logs are used for process mining, while the output files are used for KPI and disposal analysis.

The case notion is the product, represented by \texttt{productIDStr}. The activity is represented by \texttt{event}, and the timestamp is represented by \texttt{timeStamp}. Vehicle and quality information are kept as additional event or case attributes where available.

\begin{table}[t]
\caption{Main Event Log Columns}
\label{tab:event-schema}
\centering
\begin{tabular}{|l|l|}
\hline
\textbf{Column} & \textbf{Role in Analysis} \\
\hline
\texttt{productIDStr} & Case identifier \\
\hline
\texttt{event} & Activity name \\
\hline
\texttt{timeStamp} & Timestamp \\
\hline
\texttt{vehicle} & Resource identifier \\
\hline
\texttt{vehicleType} & Resource type \\
\hline
\texttt{currentDecayLevel} & Quality decay attribute \\
\hline
\texttt{processingStation} & Station attribute \\
\hline
\end{tabular}
\end{table}

\subsection{Data Processing}

The preprocessing was implemented in Python. The event logs were imported, cleaned, sorted, and transformed into a process-mining-ready format. The main preprocessing decisions were:

\begin{itemize}
\item use \texttt{productIDStr} as the product case identifier;
\item remove events without a valid product case;
\item convert simulation timestamps into a consistent timestamp format;
\item add experiment and run identifiers from file names;
\item create unique case identifiers by combining experiment, run, and product ID;
\item export cleaned data for reproducible analysis;
\item compute case-level and experiment-level KPI tables.
\end{itemize}

The analysis focuses on three selected experiments across all 20 runs each. This gives a manageable but meaningful comparison:

\begin{itemize}
\item Experiment 10: balanced fleet with random dispatching;
\item Experiment 14: balanced fleet with quality-aware and shortest-distance dispatching;
\item Experiment 23: no-UAV fleet with quality-aware and shortest-distance dispatching.
\end{itemize}

\subsection{Mining Techniques}

The project uses four main analysis techniques. First, process discovery is performed using Directly-Follows Graphs (DFGs) and an Inductive Miner model for initial validation. Second, variant analysis is used to identify dominant and rare product paths. Third, expected-lifecycle conformance analysis checks how often products follow the expected sequence from arrival to final drop-off. Fourth, enhancement and performance analysis are performed by connecting process behavior to waiting time, cycle time, quality decay, disposal, and vehicle-resource usage.

\subsection{Tools and Reproducibility}

Python was used for data preprocessing and analysis. PM4Py was used for process discovery and process mining statistics. Pandas was used for data cleaning and KPI analysis, while Matplotlib was used to generate report figures. The complete workflow is stored in a Git repository to support reproducibility.

\section{Case Description and Data}

The logistics case consists of three regions. In Region 1, products arrive and request transport. In Region 2, products are processed in a climate-controlled environment. In Region 3, products are transported to climate-controlled trucks. The fleet consists of unmanned aerial vehicles (UAVs), human-driven forklifts (HDFs), and automated guided vehicles (AGVs). Different experiments change the dispatching rules and the available fleet composition.

The expected high-level product lifecycle is:

\begin{enumerate}
\item arrival at source;
\item product calls for transport in Region 1;
\item assignment to vehicle in Region 1;
\item pickup in Region 1;
\item drop-off in Region 2;
\item start processing in Region 2;
\item finish processing in Region 2;
\item product calls for transport in Region 3;
\item assignment to vehicle in Region 3;
\item pickup in Region 3;
\item final drop-off in Region 3.
\end{enumerate}

The selected experiments provide a focused comparison between baseline behavior, quality-aware dispatching, and an alternative no-UAV fleet configuration.

\section{Analysis and Findings}

\subsection{Event Log Overview and Discovery Results}

Table~\ref{tab:log-overview} summarizes the cleaned event logs used for the main analysis. The analysis covers 60 logs in total: 3 experiments with 20 runs each.

\begin{table*}[t]
\caption{Cleaned Event Log Overview Across Selected Experiments}
\label{tab:log-overview}
\centering
\begin{tabular}{|l|p{4.2cm}|c|r|r|c|c|c|}
\hline
\textbf{Exp.} & \textbf{Configuration Role} & \textbf{Runs} & \textbf{Cases} & \textbf{Events} & \textbf{Activities} & \textbf{Variants} & \textbf{Main Variant} \\
\hline
10 & Baseline balanced fleet with random dispatching & 20 & 152752 & 1666336 & 11 & 13 & 98.35\% \\
\hline
14 & Balanced fleet with quality-aware and shortest-distance dispatching & 20 & 152678 & 1666266 & 11 & 15 & 98.44\% \\
\hline
23 & Alternative fleet without UAVs, using quality-aware and shortest-distance dispatching & 20 & 154912 & 1582452 & 12 & 17 & 85.35\% \\
\hline
\end{tabular}
\end{table*}

Experiments 10 and 14 show highly stable process behavior. In both configurations, more than 98\% of cases follow the dominant product lifecycle. Experiment 23 shows lower stability, with a main variant share of only 85.35\% and a higher number of variants. This indicates that the no-UAV configuration creates more variable process behavior.

\begin{figure*}[t]
\centering
\includegraphics[width=0.95\textwidth]{figures/Exp10_all_runs_dfg.png}
\caption{Directly-Follows Graph for Experiment 10 across all 20 runs.}
\label{fig:exp10-dfg}
\end{figure*}

\begin{figure*}[t]
\centering
\includegraphics[width=0.95\textwidth]{figures/Exp14_all_runs_dfg.png}
\caption{Directly-Follows Graph for Experiment 14 across all 20 runs.}
\label{fig:exp14-dfg}
\end{figure*}

\begin{figure*}[t]
\centering
\includegraphics[width=0.95\textwidth]{figures/Exp23_all_runs_dfg.png}
\caption{Directly-Follows Graph for Experiment 23 across all 20 runs.}
\label{fig:exp23-dfg}
\end{figure*}

The discovered DFGs show that all selected experiments follow the same general product lifecycle. However, Experiment 23 has more irregularity, which is consistent with its higher number of variants and lower dominant-variant share.

\subsection{Variant Analysis}

The first pipeline validation was performed on Experiment 10, Run 1. This log contained 7559 cases, 82507 events, 11 activities, and 10 variants. Table~\ref{tab:variants-exp10run1} shows that 98.44\% of cases in this run followed the complete expected process path.

\begin{table}[t]
\caption{Most Frequent Variants in Experiment 10 Run 1}
\label{tab:variants-exp10run1}
\centering
\begin{tabular}{|p{0.45\linewidth}|c|c|}
\hline
\textbf{Variant Description} & \textbf{Cases} & \textbf{Share} \\
\hline
Complete product flow from arrival to final drop-off & 7441 & 98.44\% \\
\hline
Starts at Region 2 completion and continues to final drop-off & 47 & 0.62\% \\
\hline
Arrival to start of processing only & 45 & 0.60\% \\
\hline
Other rare or incomplete variants & 26 & 0.34\% \\
\hline
\end{tabular}
\end{table}

The variant analysis shows that the process is mostly stable in the baseline setting. Most deviations are incomplete traces that either start late or end early. This is expected because warm-up and end-period filtering can remove some events at the boundaries of the simulation window.

\subsection{Expected Lifecycle Conformance}

Table~\ref{tab:lifecycle-conformance} shows how many cases follow the exact expected lifecycle. This analysis is used as a practical conformance-style check against the expected product sequence.

\begin{table*}[t]
\caption{Expected Lifecycle Conformance Across Selected Experiments}
\label{tab:lifecycle-conformance}
\centering
\begin{tabular}{|l|r|r|r|r|}
\hline
\textbf{Exp.} & \textbf{Exact Expected Lifecycle} & \textbf{Ends Early} & \textbf{Starts Late} & \textbf{Other Deviations} \\
\hline
10 & 98.35\% & 0.82\% & 0.82\% & 0.00\% \\
\hline
14 & 98.44\% & 0.78\% & 0.78\% & 0.00\% \\
\hline
23 & 85.35\% & 12.44\% & 2.18\% & 0.03\% \\
\hline
\end{tabular}
\end{table*}

Experiments 10 and 14 show highly conforming behavior, with 98.35\% and 98.44\% of cases following the exact expected lifecycle. Experiment 23 performs worse, with only 85.35\% exact lifecycle conformance and 12.44\% of cases ending early. This indicates that the no-UAV fleet composition creates less stable product flow and more incomplete product lifecycles.

\subsection{Performance and Quality Analysis}

Table~\ref{tab:kpi-comparison} compares the main operational and quality KPIs across the selected experiments. The values are averaged across all 20 runs per experiment.

\begin{table*}[t]
\caption{KPI Comparison Across Selected Experiments}
\label{tab:kpi-comparison}
\centering
\begin{tabular}{|l|r|r|r|r|r|r|r|}
\hline
\textbf{Exp.} & \textbf{Products} & \textbf{Wait R1} & \textbf{Wait R3} & \textbf{Total Wait} & \textbf{Travel} & \textbf{Cycle Time} & \textbf{Quality Decay} \\
\hline
10 & 134802 & 5.37 & 9.18 & 14.32 & 31.13 & 272.37 & 86.24 \\
\hline
14 & 134802 & 1.37 & 4.19 & 5.54 & 27.33 & 258.03 & 89.82 \\
\hline
23 & 128134 & 15.85 & 20.30 & 34.46 & 46.84 & 1585.85 & 42.86 \\
\hline
\end{tabular}
\end{table*}

Experiment 14 performs best in terms of waiting time and cycle time. Compared with the baseline configuration in Experiment 10, the mean total waiting time decreases from 14.32 to 5.54 time units, while the mean cycle time decreases from 272.37 to 258.03 time units. This suggests that the combination of quality-aware and shortest-distance dispatching improves operational flow without requiring a different fleet composition.

Experiment 23 shows substantially worse operational performance. Its mean total waiting time increases to 34.46 time units, and the mean cycle time reaches 1585.85 time units. This indicates that the no-UAV fleet composition creates stronger bottlenecks, especially before pickup in Region 1 and Region 3.

\begin{figure}[t]
\centering
\includegraphics[width=\linewidth]{figures/total_waiting_time_by_experiment.png}
\caption{Mean total waiting time across selected experiments.}
\label{fig:total-waiting-time}
\end{figure}

\begin{figure}[t]
\centering
\includegraphics[width=\linewidth]{figures/cycle_time_by_experiment.png}
\caption{Mean cycle time across selected experiments.}
\label{fig:cycle-time}
\end{figure}

Although Experiment 23 has a lower mean quality decay value, this should not be interpreted as better quality performance. Products with excessive quality decay may be removed from the filtered output because they are disposed earlier. Therefore, product disposal must be analyzed before drawing quality conclusions.

\subsection{Disposal Analysis}

Table~\ref{tab:disposal} shows the number and rate of disposed products. Experiments 10 and 14 do not dispose any products across the selected runs. In contrast, Experiment 23 disposes 6668 products, corresponding to a disposal rate of 4.95\%.

\begin{table}[t]
\caption{Disposed Products Across Selected Experiments}
\label{tab:disposal}
\centering
\begin{tabular}{|l|r|r|r|}
\hline
\textbf{Exp.} & \textbf{Disposed} & \textbf{Mean/Run} & \textbf{Rate} \\
\hline
10 & 0 & 0.00 & 0.00\% \\
\hline
14 & 0 & 0.00 & 0.00\% \\
\hline
23 & 6668 & 333.40 & 4.95\% \\
\hline
\end{tabular}
\end{table}

\begin{figure}[t]
\centering
\includegraphics[width=\linewidth]{figures/disposal_rate_by_experiment.png}
\caption{Product disposal rate across selected experiments.}
\label{fig:disposal-rate}
\end{figure}

The disposal analysis explains why Experiment 23 shows a lower mean quality decay in Table~\ref{tab:kpi-comparison}. Products with excessive decay are removed through disposal, so the remaining products have lower observed mean decay. Therefore, the lower quality decay value in Experiment 23 is misleading. Experiment 23 performs worse overall because it combines higher waiting time, higher cycle time, lower lifecycle conformance, and product disposal.

\subsection{Vehicle and Resource Analysis}

Table~\ref{tab:vehicle-resource} analyzes vehicle usage and mean waiting time by region. This connects the process mining results to operational resource behavior.

\begin{table*}[t]
\caption{Vehicle Usage and Mean Waiting Time by Region}
\label{tab:vehicle-resource}
\centering
\begin{tabular}{|l|l|l|r|r|r|}
\hline
\textbf{Exp.} & \textbf{Region} & \textbf{Vehicle} & \textbf{Share} & \textbf{Mean Wait} & \textbf{Mean Travel} \\
\hline
10 & Region 1 & UAV & 49.66\% & 0.93 & 6.25 \\
10 & Region 1 & HDF & 25.30\% & 5.28 & 15.00 \\
10 & Region 1 & AGV & 25.04\% & 14.27 & 30.00 \\
\hline
10 & Region 3 & UAV & 47.82\% & 1.84 & 8.33 \\
10 & Region 3 & AGV & 28.24\% & 21.11 & 37.50 \\
10 & Region 3 & HDF & 23.93\% & 9.78 & 20.00 \\
\hline
14 & Region 1 & UAV & 61.21\% & 0.41 & 6.25 \\
14 & Region 1 & HDF & 23.78\% & 1.99 & 15.00 \\
14 & Region 1 & AGV & 15.01\% & 4.32 & 30.00 \\
\hline
14 & Region 3 & UAV & 57.89\% & 1.07 & 8.33 \\
14 & Region 3 & HDF & 22.93\% & 5.90 & 20.00 \\
14 & Region 3 & AGV & 19.18\% & 11.59 & 37.50 \\
\hline
23 & Region 1 & AGV & 55.66\% & 18.61 & 30.00 \\
23 & Region 1 & HDF & 44.34\% & 12.39 & 15.00 \\
\hline
23 & Region 3 & AGV & 59.27\% & 23.02 & 37.50 \\
23 & Region 3 & HDF & 40.73\% & 16.34 & 20.00 \\
\hline
\end{tabular}
\end{table*}

UAVs are associated with the lowest waiting and travel times in both regions. Experiment 14 uses UAVs more frequently than Experiment 10, especially in Region 1, where the UAV share increases from 49.66\% to 61.21\%. This helps explain the lower waiting times observed in Experiment 14. In contrast, Experiment 23 does not use UAVs and relies only on AGVs and HDFs. This configuration produces substantially higher waiting times in both regions, especially for AGVs. The resource analysis therefore supports the recommendation to avoid the no-UAV configuration and prefer the quality-aware balanced-fleet configuration represented by Experiment 14.

\section{Improvement Plan}

The findings support three targeted process improvements.

\begin{table}[t]
\caption{Proposed Process Improvements}
\label{tab:improvements}
\centering
\begin{tabular}{|p{0.28\linewidth}|p{0.28\linewidth}|p{0.28\linewidth}|}
\hline
\textbf{Improvement} & \textbf{Expected Benefit} & \textbf{Challenge} \\
\hline
Use quality-aware and shortest-distance dispatching & Lower waiting time and cycle time & Requires real-time quality and location information \\
\hline
Avoid no-UAV fleet configuration & Prevents high waiting time and disposal & UAV availability and operating cost must be considered \\
\hline
Monitor early warning indicators for high-risk products & Reduces product disposal risk & Requires continuous event and quality tracking \\
\hline
\end{tabular}
\end{table}

The first recommendation is to use the configuration represented by Experiment 14. This configuration improves waiting time and cycle time compared with the random baseline and maintains a 0\% disposal rate. The second recommendation is to avoid the no-UAV configuration represented by Experiment 23 because it leads to more deviations, longer cycle times, and product disposal. The third recommendation is to monitor waiting time and current quality decay together so that products at risk can be prioritized before they reach the disposal threshold.

\section{Recommendations}

For the logistics stakeholder, the main recommendation is to prefer the quality-aware balanced-fleet configuration represented by Experiment 14. This configuration has the best combination of lifecycle stability, low waiting time, low cycle time, and no product disposal. The stakeholder should not select Experiment 23's no-UAV configuration, because its lower observed quality decay is caused by disposal effects and does not represent better overall performance.

For dataset and process owners, the event log quality can be improved by keeping product, vehicle, activity, timestamp, and quality decay attributes complete and consistent. The event log should also clearly distinguish product-level and vehicle-level events. This would support more advanced analysis, such as object-centric process mining of products and vehicles together.

\section{Discussion and Limitations}

This project uses simulation-generated event logs rather than real operational data. This provides controlled experiment settings, but it limits direct generalization to real-world logistics operations. The findings should therefore be interpreted as evidence within the simulated setting.

The analysis focuses on three selected experiments rather than all 27 experiments. This keeps the scope manageable and allows a focused comparison between baseline dispatching, quality-aware dispatching, and no-UAV fleet composition. However, future work could extend the analysis to all experiments.

Warm-up-filtered logs were used as the main input. This improves statistical validity, but it can also produce traces that start late or end early because some events are removed at the boundaries of the simulation window. This limitation was visible in the expected lifecycle conformance analysis.

Finally, the main case notion is the product. This is appropriate for analyzing product lifecycle and quality decay, but it does not fully capture the interaction between products and vehicles. Object-centric process mining could address this limitation in future work.

\section{Conclusion and Future Work}

This report applied process mining to a quality-aware smart-pallet logistics process. The analysis combined process discovery, variant analysis, expected lifecycle conformance, KPI analysis, disposal analysis, and vehicle-resource analysis. The results show that Experiment 14 is the strongest selected configuration. Compared with the random baseline in Experiment 10, it reduces mean total waiting time from 14.32 to 5.54 time units and reduces mean cycle time from 272.37 to 258.03 time units, while maintaining high lifecycle conformance and zero product disposal.

Experiment 23 is not recommended. Although it has a lower observed mean quality decay, this result is misleading because 6668 products are disposed, corresponding to a 4.95\% disposal rate. It also has lower expected lifecycle conformance, higher waiting time, and much higher cycle time.

Future work can extend the analysis to all 27 experiments, apply object-centric process mining to products and vehicles, and develop predictive monitoring for products at risk of excessive quality decay.

\section*{Declaration of Generative AI and AI-Assisted Technologies}

During the preparation of this report, generative AI tools were used to support brainstorming, structuring, code explanation, and language refinement. All analyses, interpretations, figures, and final decisions were reviewed and validated by the authors. The authors remain responsible for the content of the report.

\section*{Reproducibility and Artifact Availability}

All preprocessing scripts, analysis scripts, generated figures, tables, and instructions required to reproduce the results are made available in the project repository:

\begin{center}
\url{https://github.com/BoyanChakarov/process-mining-logistics}
\end{center}

The repository contains the Python scripts for loading logs, preprocessing selected experiments, generating discovery models, calculating expected lifecycle conformance, calculating KPI and disposal metrics, and analyzing vehicle-resource behavior. Due to the size of the original dataset, raw data files are not included directly in the repository. Instead, the repository contains instructions for placing the required dataset folders locally before running the scripts.

\section*{Acknowledgment}

The authors thank the course staff for guidance and the dataset creators for making the simulation data available for research and education.

\begin{thebibliography}{00}

\bibitem{van_der_aalst}
W. M. P. van der Aalst, \textit{Process Mining: Data Science in Action}, 2nd ed. Berlin, Germany: Springer, 2016.

\bibitem{pm2}
M. L. van Eck, X. Lu, S. J. J. Leemans, and W. M. P. van der Aalst, ``PM2: A process mining project methodology,'' in \textit{Advanced Information Systems Engineering}, Cham, Switzerland: Springer, 2015, pp. 297--313.

\bibitem{pm4py}
A. Berti, S. J. van Zelst, and W. M. P. van der Aalst, ``Process mining for Python: Bridging the gap between process- and data science,'' in \textit{Proc. ICPM Demo Track}, 2019.

\bibitem{bemthuis2020}
R. Bemthuis, M. Mes, M.-E. Iacob, and P. Havinga, ``Using agent-based simulation for emergent behavior detection in cyber-physical systems,'' in \textit{Proc. Winter Simulation Conference}, 2020, pp. 230--241.

\bibitem{dataset}
R. Bemthuis, M. Mes, M.-E. Iacob, and P. Havinga, ``Data underlying the paper: Using agent-based simulation for emergent behavior detection in cyber-physical systems,'' 4TU.Centre for Research Data, 2021.

\bibitem{hybrid}
N. Mustafee, A. Harper, and B. S. Onggo, ``Hybrid modelling and simulation: Driving innovation in the theory and practice of M\&S,'' in \textit{Proc. Winter Simulation Conference}, 2020, pp. 3140--3151.

\end{thebibliography}

\end{document}
